%======================================================================
\section{Discussion}\label{sec:discussion}
%======================================================================

\paragraph{Usable entanglement, not dimension.} The question we started from has a cheap version and a costly one. The cheap
version asks for a large Hilbert space; a large dimension, or a large marginal entropy, can however sit in a separable state,
and, conditionally on~\cite{Dou26}, the $I_{3322}$ example shows that near a Bell optimum dimension can diverge while usable
entanglement stays bounded. The costly version asks
whether every near-winning strategy, mixed and correlated with anything, can be made to give up Bell pairs. Theorem~A answers it
qualitatively for consistency-tested games, and the spectrum answers it quantitatively for the games we can analyse. In both,
the conclusion is drawn on the original state (\cref{sec:status}), with actual operators of the two players.

\paragraph{How the proof uses consistency.} In the classes of Theorem~A every question is checked against a copy of itself or of
a coarser question, so near-perfect play forces Alice's operators to mirror Bob's on the state. The mirrors make Bob's moments
tracial in the limit and let polynomials found in the limit act back on the actual strategy. For a game in either class with no
finite-dimensional perfect strategy, every such limit is of type~II: a type~I summand would yield a finite-dimensional perfect
strategy. A type~II algebra contains matrix algebras of every size, and this is where the Bell pairs come from. We do not know
that consistency testing is necessary for the conclusion (\cref{sec:boundary}).

\paragraph{Three mechanisms.} The quantitative results use three mechanisms, each tied to a quantitative invariant of the source. \emph{Counting}: many
approximately independent Pauli pairs force entropy and, with extraction, usable entanglement; the pairs come from lamps, and
their cost from the areas of the fillings that make lamps commute. \emph{Transport}: dilations and rotations force a spectral law
to be close to an invariant one, and invariant laws satisfy a Frostman bound of exponent $\log_mq$, hence have dimension at
least $\log_mq$; independent copies add. \emph{Addressing}: under
exact relations, the cost is the number of qubits reachable by short words. The gaps between the three are informative. Exact
and approximate relations can differ exponentially in cost, because approximate models can merge addresses or replace a sum of
defect logarithms by the logarithm of a product (\cref{prop:collective}); and the quarter exponent of the point game's upper bound
sits exactly where collective approximation beats counting.

\paragraph{What is fair to call new.} Each result is a composition. The tracial correspondences, finite von Neumann algebra
structure, Pauli rigidity and regular dilations, pretty-good recovery, uncertainty with memory and hashing are established. What
the compositions add is accounting that keeps the original state at every step, so that the conclusions are lower bounds on the
usable entanglement of every mixed near-winner of a fixed game, with explicit rates where we can compute them.

\paragraph{What remains.} Two problems are central. The matching lower exponent $\frac14$ for the point game would close the only
polynomial window we know, $[(2\log_2\lambda+1)^{-1},\frac14]$, with $\lambda$ from \cref{prop:point-fill-sharp}. Bounds by filling area cannot reach it through worst-pair raw cross-copy
estimates: that route stops at $\eps^{-1/5}$, the value of a quadratic area bound (\cref{app:window}); it would follow from a
same-state quartic bound on a Pauli multiplication table or from another original-state method. The logarithmic floor asks whether
every nonattaining consistency game demands at least logarithmic usable entanglement. It holds for every game that recovers a sign
and an even dilation with actual mirrors, including a game compiled from Abels' group (\cref{thm:floor-class,cor:abels}).
\Cref{prop:criterion} gives an operator-norm route to a counterexample; \cref{app:floor} excludes locally stable sources, the dyadic
pattern and copied Baumslag--Solitar models from that route. \Cref{thm:floor-class} separately excludes games with the stated
original-state sign-and-dilation interface. Neither the undecidability of the dichotomy nor the type~II structure of the limits
decides the problem.
